%% Chapter 14: Parameterized Algorithms I
%% Status: skeleton -- learning goals and section outline fixed, prose to be written.

\chapter{Parameterized Algorithms I}
\label{ch:fpt-i}

Parameterized complexity refuses to measure the difficulty of an instance by its
length alone. Instead we isolate a parameter $k$ --- the size of the solution,
the treewidth, the number of allowed edits --- and ask for algorithms whose
exponential behaviour is confined to $k$. Astonishingly often, this works.

\section*{Learning goals}
\addcontentsline{toc}{section}{Learning goals}

After this chapter you should be able to:
\begin{itemize}[leftmargin=*]
  \item define parameterized problems, FPT, XP and para-NP-hardness;
  \item explain the general idea of a branching algorithm in the parameterized setting;
  \item define \lang{Vertex Cover} and solve it in FPT time by branching;
  \item define \lang{Cluster Editing} and give a branching algorithm for it;
  \item give a para-NP-complete parameterized problem;
  \item define kernels and prove that a problem is FPT if and only if it has a kernel;
  \item give an $O(k^2)$ kernel for \lang{Vertex Cover};
  \item define \lang{Convex String Recoloring} and give it an $O(k^2)$ kernel;
  \item define \lang{Directed $k$-Path} and solve it in polynomial time on acyclic graphs;
  \item define constant one-sided error algorithms and give one for \lang{Directed $k$-Path} in $O(k^k n^c)$ time.
\end{itemize}

\section{Parameterized problems}
\label{sec:fpt-definitions}

\todo{Write this section.}
% Planned content: FPT, XP, para-NP-hardness, and the relations between them.

\section{Bounded search trees}
\label{sec:fpt-branching}

\todo{Write this section.}
% Planned content: \lang{Vertex Cover} in $\Ostar(2^k)$ time and \lang{Cluster Editing} in $\Ostar(3^k)$ time.

\section{Kernelization}
\label{sec:kernels}

\todo{Write this section.}
% Planned content: Definition, the equivalence of FPT and kernelizability, the $O(k^2)$ kernel for \lang{Vertex Cover}, and \lang{Convex String Recoloring}.

\section{Colour coding and random separation}
\label{sec:colour-coding}

\todo{Write this section.}
% Planned content: \lang{Directed $k$-Path}: easy on DAGs, and a one-sided-error randomized algorithm in general.

\section{Exercises}
\label{sec:ex-fpt-i}

\todo{Write this section.}

\begin{poolquestions}
  \poolq{15}{Parameterized complexity}All eleven items are covered by \cref{sec:fpt-definitions,sec:fpt-branching,sec:kernels,sec:colour-coding}.
\end{poolquestions}
